\chapter{Introdução} \label{cpt:intro}
O desenvolvimento e a pesquisa na área da saúde são temas que causam grandes impactos de forma direta na vida dos cidadãos. 
%GA: adicionar alguma ref sobre eficacia de vacinas abaixo (GV: corrigido)
Um exemplo recente foi a pandemia causada pelo contágio da COVID-19, na qual o desenvolvimento de vacinas pelos laboratórios de pesquisa minimizou a taxa de contágio do vírus, ajudou na prevenção de morbidade e mortalidade e possibilitou o retorno ao convívio social após meses de restrições~\cite{VACINAS}.

Fazer pesquisa na área da saúde não se resume a idealizar soluções, como vacinas. %apenas a nível mundial. 
%SO: Achei estranha a frase acima... parece que queremos fazer algo para além do planeta???? dei uma mexida (GV: corrigido)
Outros contextos de pesquisa também contribuem para o desenvolvimento dessa área. A pesquisa sobre melhorias na gestão de recursos públicos de saúde, por exemplo, também é uma forma de contribuir para o avanço dessa área e tem um impacto importante na população local. Desenvolver uma solução que auxilie as equipes de gestão de municípios a encontrar formas de reduzir os custos na distribuição de recursos nos hospitais da cidade pode resultar em um maior fluxo de atendimentos e até mesmo na destinação de verbas para outras áreas, como segurança e educação, por exemplo, ou ainda ajudar na prevenção da escassez de recursos~\cite{ganesh2023fair}.
%GA: refs acima? (GV: revisar)

O investimento em saúde no Brasil não tem acompanhado as demandas da população e apresentou um cenário de estagnação orçamentária na última década, crescendo apenas 2,5\%~\cite{ORCAMENTOSAUDEBR}. Em Porto Alegre, o orçamento disponibilizado para a saúde em 2024 foi de R\$ 342.149.858,00, valor que inviabiliza o atendimento mínimo das necessidades em saúde do município~\cite{ORCAMENTO}. Considerando um orçamento insuficiente e a tendência de aumento de problemas como precarização e sobrecarga, é necessária a existência de abordagens capazes de contribuir para redução de custos na área da saúde, e o uso da tecnologia pode ser um grande aliado na tentativa de diminuir essas dificuldades. 
%GA: adiciona o sobrenome dos autores quando for falar dos trabalhos deles desta forma. Ex: "Os trabalhos de ABC [14] e DED et al. [4]". Aqui também acho que caberia uma explicação do que são esses trabalhos. 1-2 frases para cada um.  (GV: corrigido)
Os trabalhos de Silva, G. F.~\cite{9239083} e El-Bouri, R. et al.~\cite{9103941} apresentam alguns exemplos de abordagens tecnológicas possíveis de serem utilizadas de modo a contribuir para o desenvolvimento da área da saúde. O primeiro utiliza a abordagem de simulação urbana virtual para fornecer \textit{insights} precisos sobre a dinâmica populacional e oferecer informações cruciais para a tomada de decisões dos gestores urbanos. Já o segundo trabalho visa facilitar a preparação dos espaços de leitos hospitalares através do treinamento de um modelo de aprendizado profundo.

Observando que o orçamento para investimentos em saúde em Porto Alegre é limitado e insuficiente, espera-se que as soluções propostas não tenham um custo elevado para sua implantação. Com isso, simular ambientes urbanos virtuais pode ser uma alternativa de baixo custo a ser explorada na tentativa de redução de custos em saúde.

Ao trabalhar com a simulação de ambientes urbanos virtuais, podemos estudar o comportamento de cenários em diversas situações. A superlotação de estádios de futebol, a transmissão de uma doença contagiosa e o trânsito de pessoas pelos bairros da cidade são alguns exemplos de aplicações. Os resultados dessas simulações podem ser objetos de estudo para prevenir eventos prejudiciais à população, realizar ações de marketing e até mesmo entender o comportamento de eventos ocorridos no passado. São diversas as possibilidades de aplicações. Na área da saúde, simular grupos de pessoas em cidades com muitos habitantes, agregar dados da vida para um ambiente virtual e gerar dados para estudos são alguns dos desafios associados a essa área. Uma aplicação da simulação de um ambiente urbano virtual nesse contexto pode ser empregado para estimar o fluxo de pacientes nos leitos de um ou mais hospitais nas cidades durante um período parametrizado. Essa dissertação visa a proposta de desenvolvimento do LODUS-Health, um simulador que trata de alguns temas da saúde em um ambiente de simulação urbana virtual.

%SO: Aqui falta dizer o desafio da área... p.ex. desafio de simular pessoa por pessoa em cidades com miutos habitantes, necessidade de agregar dados da vida real... e etc... No final dizer que essa dissertação visa desenvolver um simulador que permita tartar de alguns temas da saúde dentro da simulação urbana. (GV: corrigido)

\section{Objetivo}
%GA: para framework e outras palavras estrangeiras, usar itálico (GV: corrigido)
O objetivo deste trabalho é apresentar a pesquisa e desenvolvimento do LODUS-Health, uma ferramenta de simulação do fluxo de ocupação dos leitos hospitalares da cidade de Porto Alegre, 
%SO: Vaz: tu achas que tu tu fizeste um simulador de fluxod e ocupação de leitos, ou seria mais adequado dizer que tu simulas fluxos de pessoas na cidade que se direcionam a ociupar hospitais? ou algo assim (GV: revisar)
visando aprimorar a gestão hospitalar pública em diversos cenários, como na destinação de recursos. Espera-se que a ferramenta seja capaz de prover aos usuários a capacidade de realizar simulações de um ambiente configurado com dados reais de quantidade de leitos hospitalares e permita que usuários especialistas na área da saúde ou gestores urbanos ou hospitalares possam observar o comportamento dos leitos hospitalares em um contexto customizado. Para isso, é desenvolvida uma abordagem de simulação com configurações parametrizáveis utilizando o \textit{framework} de simulação urbana virtual LODUS~\cite{9239083}, geração de textos médicos de evolução de pacientes com a OpenAI API\footnote{\url{https://openai.com/blog/openai-api}} e uma interface gráfica com ferramentas de desenvolvimento web. Conforme será descrito nos resultados, o modelo proposto neste projeto permite a simulação de cenários de ocupação de leitos em diferentes especialidades dentro dos hospitais de Porto Alegre. %GA: ref e footnote para o LODUS e para a API da OpenAI (GV: revisar)

%SO: para ter uma section só para objetivos ela deve incluir objetivos específicos. Exemplos: integrar com dados reais da cidade, como doenças CIDI e dados de hopitais, prover ferramenta ao usuários final que permita blabla.... (GV: corrigido)


%SO: Ao final, uma frase, p.ex.: " Conforme será descrito nos resutlados, o modelo proposto neste projeto permite o desenvolvimento de cenários para simulação urbana e de encaminhamento de doentes aos hospitais. (GV: revisar)

\section{Organização do Trabalho}
Esta pesquisa está organizada nos seguintes capítulos: O Capítulo~\ref{cpt:contexto} contextualiza sobre as principais tecnologias utilizadas no desenvolvimento do modelo de simulação, são descritos o \textit{framework} LODUS e a OpenAI API nas seções \ref{sec:back_lodus} e \ref{sec:chat}. No Capítulo~\ref{cpt:lit}, é apresentada a metodologia utilizada na seleção dos trabalhos relacionados, descrevendo a fonte de busca (Seção~\ref{sec:fontes_busca}), os critérios de inclusão e exclusão (Seção~\ref{sec:criterios}), o processo de seleção de materiais (Seção~\ref{sec:processo_selecao}), os trabalhos relacionados (Seção~\ref{sec:trabalhos_selecao}) e uma breve discussão sobre eles é feita na Seção~\ref{sec:trabalhos_discussao}. O Capítulo~\ref{cpt:met_prop} discorre sobre a metodologia aplicada no desenvolvimento da pesquisa, passando pelas etapas de criação do conjunto de dados utilizado, incluindo a coleta de dados (Seção~\ref{sec:sites}), a estruturação do conjunto de dados e uma análise sobre eles (Seção~\ref{sec:montagem_conjunto}). Também é apresentada as etapas de desenvolvimento da interface gráfica (Seção~\ref{sec:interface_config}), do \textit{plugin} de simulação (Seção~\ref{sec:plugin_lodus}), a visualização dos resultados do modelo de simulação (Seção~\ref{sec:visualizacao_resultados}) e geração de texto médico individualizado (Seção~\ref{sec:geracao_texto}). O Capítulo~\ref{cpt:result} apresenta os resultados da simulação em quatro diferentes contextos. Por fim, os Capítulos~\ref{cpt:discus} e~\ref{cpt:conclus} apresentam uma discussão sobre o trabalho e a conclusão, respectivamente.

%SO: de maneira geral, a section precisa de mais refs (GV: corrigido)